% ============================================================
%  Kapitel 1 — Einleitung
% ============================================================
\section{Einleitung}

Supraleitung ist eines der faszinierendsten Phänomene der Festkörperphysik: Unterhalb einer materialspezifischen kritischen Temperatur $T_c$ verschwindet der elektrische Widerstand eines Leiters vollständig, und das Material expelliert das Magnetfeld aus seinem Inneren
(Meißner-Ochsenfeld-Effekt)~\cite{Meissner1933}.
Seit der Entdeckung der Supraleitung in Quecksilber durch Heike Kamerlingh Onnes im Jahr 1911 hat das Phänomen die Grundlagenforschung sowie zahlreiche technologische Anwendungen~--- von Kernspin\-tomographen über Teilchenbeschleuniger bis zu supraleitenden Quantencomputern~--- maßgeblich geprägt.

Das mikroskopische Verständnis gelang erst 1957 durch Bardeen, Cooper und Schrieffer
(BCS-Theorie~\cite{BCS1957}): Elektronen bilden unterhalb von $T_c$ über phononvermittelte Wechselwirkung Cooperpaare mit Gesamtimpuls null; diese kondensieren in einen kollektiven Grundzustand, der elektrischen Widerstand verhindert.
Eine phänomenologische Beschreibung des Magnetfeldverhaltens liefern die London-Gleichungen~\cite{London1935}, die unter anderem den Meißner-Effekt quantitativ erklären.
Eine Erweiterung auf temperaturabhängige Ordnungsparameter bietet die Ginzburg-Landau-Theorie, die den Übergang in supraleitende Dünne-Film-Geometrien besonders zugänglich macht.

Mit der Entdeckung der Hochtemperatur-Supraleitung in Kuprat-Verbindungen durch Bednorz und Müller~\cite{Bednorz1986} erweiterte sich das Anwendungsgebiet erheblich: Die Sprungtemperaturen liegen weit über dem Siedepunkt von Stickstoff ($77~\mathrm{K}$), was kostengünstige Kryotechnik ermöglicht.

Im vorliegenden F-Praktikumversuch werden grundlegende Eigenschaften der klassischen Typ-I-Supraleitung an einem dünnen Bleifilm ($\mathrm{Pb}$, $100~\mathrm{nm}$) bei Temperaturen nahe dem Siedepunkt von flüssigem Helium ($T \approx 4{,}2~\mathrm{K}$) untersucht.
Konkret werden das magnetische Phasendiagramm $B_c(T)$, das die Grenze zwischen supraleitendem und normalem Zustand beschreibt, sowie der kritische Strom $I_c(T)$ als Funktion der Temperatur gemessen.
Beide Größen folgen innerhalb der Ginzburg-Landau-Theorie charakteristischen Potenzgesetzen, die an die Messdaten angepasst werden.
Ergänzend wird das Levitationsphänomen an einem YBCO-Hochtemperatursupraleiter qualitativ beobachtet.

\medskip\noindent
\textbf{Gliederung.}
Kapitel~\ref{sec:theorie} legt die theoretischen Grundlagen der Typ-I- und Typ-II-Supraleitung dar.
Der Versuchsaufbau wird in Kapitel~\ref{sec:aufbau} beschrieben.
Die Auswertung der $B_c(T)$- und $I_c(T)$-Messungen erfolgt in Kapitel~\ref{sec:auswertung}, gefolgt von Diskussion und Fazit in Kapitel~\ref{sec:diskussion}.
